\section*{第 \ref{ch:b1:organomagnesium} 章——有机镁试剂}

\begin{solution}{exo:b1:organomagnesium:1}
\ce{CH3CH2Br + Mg -> CH3CH2MgBr}；\ce{C6H5Br + Mg -> C6H5MgBr}（于无水
乙醚中）。C--Br：Br（2.96）的\omterm{def:b1:periodicity:electronegativity}{电负性}大于 C（2.55），碳
带正电。C--Mg：Mg（1.31）的\omterm{def:b1:periodicity:electronegativity}{电负性}较小，碳带负电。
\end{solution}

\begin{solution}{exo:b1:organomagnesium:2}
水：\ce{CH3MgI + H2O -> CH4 + Mg(OH)I}。
乙醇：\ce{CH3MgI + CH3CH2OH -> CH4 + CH3CH2OMgI}。
乙酸：\ce{CH3MgI + CH3COOH -> CH4 + CH3COOMgI}。
重水：\ce{CH3MgI + D2O -> CH3D + Mg(OD)I}。
\end{solution}

\begin{solution}{exo:b1:organomagnesium:3}
甲醛：丙-1-醇（伯醇）。乙醛：丁-2-醇（仲醇）。
丙酮：2-甲基丁-2-醇（叔醇）。苯甲醛：
1-苯基丙-1-醇（仲醇）。二氧化碳：丙酸。
\end{solution}

\begin{solution}{exo:b1:organomagnesium:4}
其羟基的酸性足够强，一旦生成便会破坏试剂：
分子会使自身（或其邻近分子）质子化，得到
丁-1-醇的镁醇盐 \ce{CH3(CH2)3OMgBr}，并非可用的
试剂。羟基必须先加以保护
（见\cref{ch:b1:acetals-protection}，以醛为例说明了该思路）。
\end{solution}

\begin{solution}{exo:b1:organomagnesium:5}
C--Mg 键的电子对进攻乙醛的羰基碳，C=O 的 $\pi$ 电子对
转移到 O 上，生成 \ce{CH3CH(CH3)O^-} 与 \ce{MgBr+}；随后 \ce{H3O+} 将
质子交给氧，得丙-2-醇。丙-2-醇的羟基碳上连有两个甲基，
故为非\omterm{def:b1:stereochemistry-in-depth:stereocentre}{手性}分子。（若 R 不同，例如
溴化乙基镁与乙醛反应，产物丁-2-醇具有\omterm{def:b1:stereochemistry-in-depth:stereocentre}{手性}，但
为外消旋体：平面羰基两面被等概率进攻。）
\end{solution}

\begin{solution}{exo:b1:organomagnesium:6}
羟基碳上连有一个甲基和两个乙基：
溴化乙基镁与丁-2-酮反应，或溴化甲基镁与
戊-3-酮反应。
\end{solution}

\begin{solution}{exo:b1:organomagnesium:7}
先以 2-溴丙烷与镁在无水乙醚中制备 \ce{(CH3)2CHMgBr}，
将其倾于固体 \ce{CO2} 上，再酸化得 \ce{(CH3)2CHCOOH}。$n = 12.3/123.0 = 0.100$ 摩尔；
$0.100 \times 0.70 \times 88.0 = \qty{6.2}{g}$。
\end{solution}

\begin{solution}{exo:b1:organomagnesium:8}
$0.18/18.0 = \qty{0.010}{mol}$ 的水破坏了 \qty{0.010}{mol} 的
试剂：\qty{20}{\percent}；生成苯，\ce{C6H5MgBr + H2O -> C6H6 +
Mg(OH)Br}。
\end{solution}

\begin{solution}{exo:b1:organomagnesium:9}
该试剂为\omterm{def:b1:electronic-effects:nucleophile}{亲核试剂}，与仍存在的溴苯反应：
总反应为 \ce{C6H5MgBr + C6H5Br -> C6H5-C6H5 + MgBr2}。将溴苯
缓慢加入过量镁中，可使其浓度保持较低：它遇到的是新鲜的
金属，而非已生成的试剂。
\end{solution}

\begin{solution}{exo:b1:organomagnesium:10}
1-苯基乙-1-醇：\ce{CH3MgBr}（由溴甲烷制得）与苯甲醛反应，
或 \ce{C6H5MgBr}（由溴苯制得）与乙醛反应。2-苯基丙-2-醇：
\ce{C6H5MgBr} 与丙酮反应，或 \ce{CH3MgBr} 与
1-苯基乙-1-酮反应。
\end{solution}

\begin{solution}{exo:b1:organomagnesium:11}
将酮加入试剂中，使试剂始终留在制备它的
无水烧瓶中，每份加料放出的热量被
沸腾的乙醚吸收。单用水会得到胶状的
碱式盐 \ce{Mg(OH)Br}，它会包裹产物；冷的稀酸
将其溶解为 \ce{Mg^2+}，低温还可减少
叔醇在酸中发生的副反应。
\end{solution}

\begin{solution}{exo:b1:organomagnesium:12}
每摩尔试剂与水反应生成一摩尔碱式镁盐，
\ce{RMgBr + H2O -> RH + Mg(OH)Br}，其 \ce{OH-} 被酸\omterm{def:b1:titration-methods:titration}{滴定}：
试剂的物质的量等于所用酸的量，即 \qty{0.088}{mol}，\omterm{def:b1:extent-q-and-k:fractional-extent}{产率}为
\qty{88}{\percent}。
\end{solution}

\begin{solution}{pb:b1:organomagnesium:1}
\textbf{1.} \ce{C6H5Br + Mg -> C6H5MgBr}。
\textbf{2.} Mg $0.535/24.3 = \qty{0.0220}{mol}$；\ce{C6H5Br} $3.14/156.9 =
\qty{0.0200}{mol}$：镁稍过量。
\textbf{3.} 水破坏试剂：\ce{C6H5MgBr + H2O -> C6H6 +
Mg(OH)Br}。
\textbf{4.} 阻止了空气中的水蒸气经
冷凝器进入。
\textbf{5.} 乙醚向镁提供\omterm{def:b1:lewis-resonance-vsepr:lewis}{孤电子对}（\omterm{def:b1:lewis-resonance-vsepr:lewis-acid}{路易斯碱}）：
若无此作用，试剂既不能生成，也不能留在溶液中。
\textbf{6.} \ce{C6H5MgBr + C6H5Br -> C6H5-C6H5 + MgBr2}；联苯
$0.15/154.0 = \qty{9.7e-4}{mol}$，消耗了 \qty{9.7e-4}{mol} 的
溴苯及等量的试剂：合计损失溴苯
\qty{1.9e-3}{mol}。
\textbf{7.} $3.28/182.0 = \qty{0.0180}{mol}$。
\textbf{8.} C--Mg 键的电子对进攻羰基碳，$\pi$ 电子对
转移到氧上：\ce{(C6H5)3C-O^-} \ce{MgBr+}。
\textbf{9.} 可用试剂至多 $0.0200 - 0.0019 = \qty{0.0181}{mol}$；
二苯甲酮 \qty{0.0180}{mol}：二苯甲酮刚好为\omterm{def:b1:extent-q-and-k:max-extent}{限量反应物}。
\textbf{10.} 否：羟基碳上连有三个相同的苯基。
\textbf{11.} 试剂留在其无水烧瓶中，每份酮的热量
被沸腾的乙醚吸收。
\textbf{12.} 会破坏部分试剂（生成苯），
并使二苯甲酮未参与反应。
\textbf{13.} \ce{(C6H5)3COMgBr + H3O+ -> (C6H5)3COH + Mg^2+ + Br- + H2O}。
\textbf{14.} 酸能溶解单独用水时会沉淀的
碱式镁盐；低温可抑制
第 18 题所述\omterm{def:b1:electronic-effects:intermediates}{碳正离子}的生成。
\textbf{15.} 三苯甲醇和联苯在乙醚层中；
镁盐在水层中。
\textbf{16.} 用冷石油醚洗涤（或研磨）固体：
联苯溶解，三苯甲醇则留存。
\textbf{17.} 熔点明锐，与表中数值相同（或
薄层色谱法中呈单一斑点）。
\textbf{18.} \ce{(C6H5)3COH + H+ -> (C6H5)3C+ + H2O}：叔\omterm{def:b1:electronic-effects:intermediates}{碳正离子}，
其电荷通过共振分散于三个苯环上，
非常稳定，并因\omterm{def:b1:structure-spectroscopy:conjugation}{共轭体系}扩大而呈色。
\textbf{19.} $3.75/260.0 = \qty{0.0144}{mol}$。
\textbf{20.} $0.0144/0.0180 = \qty{80}{\percent}$。
\textbf{21.} 转移及重结晶过程中的损失；试剂
被微量水分破坏；反应不完全。
\textbf{22.} 苯甲酸 \ce{C6H5COOH}，经酸化后得到。
\textbf{23.} $0.0181 \times 122.0 = \qty{2.21}{g}$。
\textbf{24.} \ce{C6H5MgBr + CO2 -> C6H5COOMgBr}，再
\ce{C6H5COOMgBr + H3O+ -> C6H5COOH + Mg^2+ + Br- + H2O}。
\textbf{25.} \textbf{\qty{80}{\percent}}。
\end{solution}
